\chapter{Giardiasis}

\section*{Definition and Overview}
Giardiasis is an infection of the small intestine caused by the microscopic parasite \textit{Giardia lamblia} (also called \textit{Giardia intestinalis} or \textit{duodenalis}). It is one of the most common parasitic causes of diarrhoeal illness worldwide. The parasite is spread through contaminated water, food, or surfaces, and from person to person, particularly in settings such as childcare centres, camps, and among travellers. Many people infected with giardia have no symptoms, while others develop diarrhoea, abdominal cramps, bloating, nausea, and foul-smelling stools that can persist for weeks. Giardiasis is treatable with specific medicines, and most people recover fully, although symptoms can sometimes linger.

\section*{Epidemiology}
Giardiasis occurs throughout the world and is a major cause of waterborne disease. In Australia, it is one of the more commonly notified infections, with several thousand cases reported each year. High-risk groups include children in childcare, their families and carers, campers and hikers who drink untreated water, travellers to countries with poor sanitation, men who have sex with men, and people with weakened immune systems. Outbreaks are often linked to contaminated drinking water supplies, swimming pools, and contaminated food. Because many infections are asymptomatic, the true number of cases is higher than reported.

\section*{Aetiology and Pathophysiology}
Giardiasis is caused by the parasite \textit{Giardia lamblia}, which exists in two forms: a hardy cyst that survives outside the body and a motile trophozoite that causes disease in the intestine.

\begin{itemize}
    \item \textbf{Cysts:} the infectious form, passed in the faeces of infected people and animals; they can survive for months in cold water and are resistant to chlorination at usual levels
    \item \textbf{Transmission:} swallowing contaminated water (from streams, lakes, or inadequately treated supplies), contaminated food, or faecal material via hands, or during sexual contact
    \item \textbf{Ingestion:} cysts travel to the small intestine, where they release trophozoites
    \item \textbf{Trophozoites:} attach to the lining of the small intestine with a suction disc, interfering with absorption and causing inflammation
    \item \textbf{Shedding:} trophozoites form new cysts that are passed in stool, continuing the cycle
    \item \textbf{Low infective dose:} as few as ten cysts can cause infection
\end{itemize}

\section*{Clinical Features}
Symptoms usually begin one to three weeks after exposure and can range from none at all to severe illness.

\begin{itemize}
    \item Watery or loose diarrhoea that may be explosive and foul-smelling
    \item Abdominal cramps, bloating, and excess wind
    \item Nausea, loss of appetite, and weight loss
    \item Stools that are greasy and float because of fat malabsorption
    \item Fatigue and a general feeling of being unwell
    \item Dehydration, especially in children and older people
    \item In some people, lactose intolerance that persists for weeks after the infection clears
    \item Asymptomatic carriage, which can still spread the parasite to others
\end{itemize}

\section*{Differential Diagnosis}
Many infections and conditions cause similar symptoms, so diagnosis needs to be confirmed.

\begin{itemize}
    \item Other intestinal parasites, including cryptosporidium and amoebiasis
    \item Bacterial gastroenteritis caused by campylobacter, salmonella, shigella, and \textit{E. coli}
    \item Viral gastroenteritis, including norovirus and rotavirus
    \item Travellers' diarrhoea from other causes
    \item Irritable bowel syndrome, which can follow an episode of gastroenteritis
    \item Coeliac disease and other causes of malabsorption
    \item Inflammatory bowel disease
    \item Food intolerance or medicine-related diarrhoea
\end{itemize}

\section*{When to Seek Medical Care}
See a doctor if diarrhoea persists beyond a few days, is severe, or is accompanied by dehydration, especially in young children, older people, or people with weakened immunity. People who have been in a high-risk situation, such as drinking untreated water while camping, should mention this to their doctor.

\textbf{Call 000 (or local emergency services) for:}
\begin{itemize}
    \item Signs of severe dehydration: very little urine, extreme thirst, sunken eyes, or drowsiness
    \item Blood in the stool or severe abdominal pain
    \item Diarrhoea in a baby under six months, especially with vomiting and poor feeding
    \item A person who is unable to keep fluids down for more than 24 hours
\end{itemize}

\section*{Diagnosis and Investigations}
Laboratory testing is needed to confirm giardiasis because the symptoms are not specific.

\begin{itemize}
    \item \textbf{Stool microscopy:} examination of stool samples for cysts and trophozoites; several samples on different days may be needed because shedding is intermittent
    \item \textbf{Antigen tests:} rapid stool tests detect giardia proteins and are more sensitive than microscopy
    \item \textbf{PCR:} molecular testing of stool can detect the parasite's DNA
    \item \textbf{Endoscopy:} rarely, small-bowel biopsy is needed when the diagnosis is elusive in persistent cases
    \item \textbf{Contact and exposure history:} asking about travel, water sources, childcare, and contacts guides testing
\end{itemize}

\section*{Management}
Treatment is effective and usually simple, though some people need a second course.

\begin{itemize}
    \item \textbf{Antiparasitic medicines:} metronidazole or tinidazole are the usual treatments, taken for five to seven days
    \item \textbf{Nitazoxanide:} an alternative for children and some adults
    \item \textbf{Fluid replacement:} oral rehydration solution or water with electrolytes to replace lost fluids
    \item \textbf{Diet:} bland, easily digested foods, and temporary avoidance of dairy and fatty foods if they aggravate symptoms
    \item \textbf{Treating contacts:} household members with symptoms should be tested and treated
    \item \textbf{Immunocompromised people:} may need longer or repeated treatment and specialist care
    \item \textbf{Notification:} giardiasis is notifiable in Australia, so cases are reported to public health authorities
\end{itemize}

\section*{Complications}
\begin{itemize}
    \item Prolonged diarrhoea and weight loss lasting weeks to months
    \item Dehydration and electrolyte disturbances, especially in children and older people
    \item Malabsorption of fats, lactose, and vitamins, leading to deficiency
    \item Failure to thrive in young children
    \item Post-infectious irritable bowel syndrome
    \item Chronic infection and persistent symptoms in people with weakened immunity
    \item Reactive arthritis, which rarely follows infection
\end{itemize}

\section*{Prognosis and Outcomes}
With treatment, most people with giardiasis recover completely within one to two weeks, and even untreated the infection usually resolves within a few weeks. A minority experience persistent symptoms, including intermittent diarrhoea, bloating, and fatigue, which can last for months and sometimes lead to a diagnosis of post-infectious irritable bowel syndrome. In people with normal immunity, complications are uncommon and the parasite is cleared effectively. In immunocompromised people, the infection may be harder to eradicate and requires specialist management.

\section*{Prevention and Self-Care}
\begin{itemize}
    \item Practise good hand hygiene, especially after using the toilet, changing nappies, and before preparing food
    \item Boil, filter, or otherwise treat water from streams, lakes, and rivers before drinking it while camping or hiking
    \item Avoid swallowing water when swimming in pools, lakes, and rivers
    \item Wash fruit and vegetables thoroughly, and wash hands after handling raw food
    \item In childcare settings, maintain strict hygiene and exclude children with diarrhoea until cleared
    \item Do not share towels, flannels, or eating utensils with an infected person
    \item Drink plenty of fluids if you become infected, and avoid dairy and fatty foods during the illness
    \item Wash your hands after contact with pets, as animals can carry giardia
\end{itemize}

\section*{Sources and Further Reading}
The following reputable sources provide further information for healthcare students and patients seeking reliable, current advice about giardiasis.

\begin{itemize}
    \item Healthdirect Australia. \textit{Giardiasis.} https://www.healthdirect.gov.au/
    \item Australian Government Department of Health and Aged Care. \textit{Giardiasis fact sheet.} https://www.health.gov.au/
    \item NSW Health. \textit{Giardiasis fact sheet.} https://www.health.nsw.gov.au/
    \item Centers for Disease Control and Prevention. \textit{Giardia fact sheet.} https://www.cdc.gov/
    \item World Health Organization. \textit{Giardiasis.} https://www.who.int/
    \item Mayo Clinic. \textit{Giardia infection (giardiasis).} https://www.mayoclinic.org/
\end{itemize}
